% Experimental phases at a glance.
% V6: wording pass; the Phase B2 / C2 protocol is now described precisely
% (saturation FIXED at 0 = full grayscale, additive noise off) -- the V5
% phrase "minus the saturation axis" wrongly suggested color was kept.
\section{Experimental phases at a glance}
\label{sec:phases}

The campaign breaks the THz robustness question into seven phases. Each phase
answers one question, and every later section refers back to this overview
for the pipeline definitions instead of re-explaining them.
Table~\ref{tab:phases_overview} is the single point of reference.

\textbf{Phase A (Clean baseline).}\quad Pipeline: identity (upsample only).
Six cells (three models $\times$ two datasets). Establishes the upper-bound
accuracy that every later phase is measured against.

\textbf{Phase B1 (Combined degradation).}\quad Pipeline: saturation lerp
$\rightarrow$ downsample $\rightarrow$ noise and salt-and-pepper at the low
resolution $\rightarrow$ upsample to $224 \times 224$ $\rightarrow$ blur,
with all five axes active at the same level. Thirty cells (three models
$\times$ two datasets $\times$ five levels). Question: how bad is the
realistic worst case, where everything degrades at once?

\textbf{Phase C (Single-axis isolation).}\quad Pipeline: exactly one axis
active at level $L$, the other four held at identity. One hundred fifty cells
(three models $\times$ two datasets $\times$ five levels $\times$ five axes).
Question: which individual degradation drives the Phase~B1 collapse?

\textbf{Phase D (Regularization-recovery sweep).}\quad Pipeline: identical to
Phase~B1, plus one of three training-time regularization treatments T1 / T2 /
T3 (architectural dropout, label mixing, and their combination). Ninety cells.
Question: is the Phase~B1 collapse over-fitting --- which regularization
should fix --- or lost information, which it cannot fix?

\textbf{Phase B2 (THz-protocol simplification).}\quad Pipeline: Phase~B1 with
the image held at full grayscale (saturation fixed at 0) and the additive
Gaussian noise turned off; T3 regularization active. Thirty cells. This is
the deployment-realistic protocol: a real THz frame is a single-channel
image, so the color axis is not a variable, and the dominant sensor noise is
impulsive rather than Gaussian. Question: how much easier is the realistic
protocol than the worst-case one?

\textbf{Phase B2-nr (No-regularization arm).}\quad Same pipeline as
Phase~B2, at L3 only, with T3 off. Six cells. Question: of the Phase
B1$\rightarrow$B2 improvement, how much comes from the simpler protocol
itself and how much from the regularization?

\textbf{Phase C2 (Single-axis isolation under the THz protocol).}\quad Same
identity-elsewhere logic as Phase~C, restricted to the three THz-relevant
axes (resolution, blur, salt-and-pepper), with grayscale and zero Gaussian
noise forced throughout. Ninety cells. Question: under realistic conditions,
which degradation axis should a THz system engineer fight first? The
Section~\ref{sec:tests} plots answer this on the measured PSNR / SSIM scales,
never on the preset L1--L5 index.

% Table II: Phases at a glance (NEW per PRD V2 \S 7.4 R4).
% Columns: phase, cells, axes active at level L, levels swept, regularization,
% primary research question. One row per phase + one row for the multi-seed
% audit.
\begin{table}[htbp]
  \centering
  \caption{Experimental phases at a glance. The campaign comprises 402
    canonical training cells at seed 42 plus 48 multi-seed audit replicates
    at seeds 43 and 44. Each phase isolates a different facet of the
    accuracy collapse under the five-level THz-like degradation curve.}
  \label{tab:phases_overview}
  \small
  \setlength{\tabcolsep}{4pt}
  \renewcommand{\arraystretch}{1.15}
  \begin{tabular}{l c p{2.6cm} p{1.6cm} p{1.4cm} p{5.3cm}}
    \toprule
    \textbf{Phase} & \textbf{Cells} & \textbf{Axes active at L} & \textbf{Levels swept} & \textbf{Reg.} & \textbf{Primary research question} \\
    \midrule
    A      & 6   & none (identity) & none & none & Upper-bound clean accuracy per (model, dataset). \\
    B1     & 30  & all 5 axes      & L1--L5 & none & Magnitude of the realistic combined-degradation collapse. \\
    C      & 150 & one axis (4 idle) & L1--L5 & none & Per-axis decomposition of the Phase B1 collapse. \\
    D      & 90  & all 5 axes      & L1--L5 & T1 / T2 / T3 & Is the Phase B1 collapse overfit or information bottleneck? \\
    B2     & 30  & 3 axes (sat $=0$, noise $=0$) & L1--L5 & T3 & Effect of the realistic THz protocol on the combined curve. \\
    B2-nr  & 6   & same as B2        & L3 only & none & Decomposition of B1$\to$B2 into protocol versus regularization. \\
    C2     & 90  & one of \{res, blur, S\&P\} & L1--L5 & T3 & Per-axis attribution under the THz protocol. \\
    \midrule
    Multi-seed audit & 48 & all 24 L3 headline cells & L3 only & inherits cell & Run-to-run variance bound on every headline cell. \\
    \bottomrule
  \end{tabular}
\end{table}

